% Case 2: Let's Talk TV. Numbers come from sections/numbers-lttv.tex
% (generated by scripts/make_tex_numbers.py); quotations from the files named
% in notes/source-verification.md. Draft, 2026-10-03.

\section{Unbundling Canadian television}
\label{sec:lttv}

\subsection{The decision and the forecast}

In March 2015 the CRTC told television distributors to offer, by March 2016, an entry-level service \enquote{priced at no more than \$25 (not including equipment) per month} \citep[para.~26]{crtc2015worldofchoice}. Every discretionary channel was to be available on its own or in small packages by March 2016, and in both forms by December 2016 \citep[para.~47]{crtc2015worldofchoice}. The policy came out of the \textit{Let's Talk TV} proceeding, at which, by one later account, \enquote{virtually all parties expressed concerns about potential impacts of unbundling} \citep[para.~xxii, p.~5]{nordicity2015television}.

Nine months later Nordicity and Peter Miller published \textit{Canadian Television 2020}, prepared for ACTRA, the Canadian Media Guild, the Directors Guild of Canada, Friends of Canadian Broadcasting and Unifor \citep{nordicity2015television}. It modelled two worlds from 2015 to 2020: a baseline with \enquote{technological and consumer behaviour changes, assuming that the pre-existing broadcasting regulatory regime remains}, and the same world with the CRTC's decisions overlaid \citep[para.~ii, p.~1]{nordicity2015television}. Every headline number is the gap between the two.

By 2020, the report concluded, the decisions were \enquote{likely to result in a loss of \lttvJobsTotal{} FTEs of employment in the Canadian economy}, along with \$\lttvGDPTotal{} million of GDP. Of those jobs, \lttvJobsDirect{} were in broadcasting and production and \lttvJobsSpinoff{} in other sectors \citep[para.~xliv and Table~1, p.~12]{nordicity2015television}. Those figures sit at the end of a chain. Subscribers move to cheaper unbundled packages; distributors lose retail revenue and pass most of the loss to Canadian channels as lower wholesale fees; weaker channels close; spending on Canadian programming falls; and input-output multipliers turn the spending into jobs and GDP.

The report states the assumptions that drive each link. Subscribers taking the entry-level service with a few chosen channels make up \lttvByopSixteen{} of all subscribers in 2016, \lttvByopSeventeen{} in 2017 and \lttvByopEighteen{} from 2018 \citep[Table~18, p.~76]{nordicity2015television}. Distributors pass \enquote{75\% of their retail revenue loss to Canadian programming services} \citep[para.~207, p.~77]{nordicity2015television}. Ten per cent of the large integrated companies' Category A and B channels and a quarter of the independents' close by 2020 \citep[paras.~232--233, p.~85]{nordicity2015television}.

By 2020 the model takes \$\lttvSpecImpactTwenty{} million (\lttvSpecShareTwenty{}) a year out of specialty and pay channels' revenue and \$\lttvBDUImpactTwenty{} million (\lttvBDUShareTwenty{}) out of distributors' revenue, against the baseline \citep[paras.~xxvi--xxvii]{nordicity2015television}.

\subsection{What the Heritage Committee heard}

On 12 April 2016 Ian Morrison of Friends of Canadian Broadcasting, accompanied by Miller, put the report to the Commons Standing Committee on Canadian Heritage. Its key findings, he said, were that \enquote{by 2020, some 15,130 media jobs will be lost}, along with a fall in Canadian programming expenditure that was \enquote{18\% of what now exists}, \enquote{all of this as the direct result of Let's Talk TV regulatory changes} \citep[time marks 0900--0905]{commons2016chpc8}.

Each phrase is stronger than the report. The \lttvJobsTotal{} jobs are economy-wide, and \lttvJobsSpinoff{} of them are spin-off employment in other sectors; the \lttvCPEShareStated{} is a share of the report's projected 2020 baseline, not of current spending; and the report itself splits its total into direct and spin-off impact. The figure also fails to reproduce. The report's own baseline-scenario figures put 2020 Canadian programming expenditure at \$\lttvCPEBaselineTwenty{} million, so its \$\lttvCPEImpactTwenty{} million reduction is \lttvCPEShareTwenty{} of baseline, not the \lttvCPEShareStated{} its text states \citep[para.~239; Figs.~8, 9, 43]{nordicity2015television}. I tested eleven candidate denominators built from the report's components; none gives \lttvCPEShareStated.

Morrison also offered a forecast of his own, one month after the entry-level service launched: \enquote{the best estimates I've seen, Mr.~Waugh, are that about 4\% of Canadians will go for it} \citep[time marks 0920--0925]{commons2016chpc8}.

\subsection{How the forecast was checked}

The prediction record (forecasters, quantities, paths by year, stated assumptions) was committed before any outcome series was downloaded. I then re-read the report's two paths from the data labels on its charts, checking every value against the totals in its text, and ran a design analysis by fake-data simulation before opening any CRTC file \citep{gelman2014types, gelman2020workflow}. It asked, for each forecast quantity, whether the outcome could separate a world in which the report's impact occurred from one in which it didn't, given that the baseline itself could be wrong.

It depended on the quantity. With the baseline's error set at the volatility of the report's own 2010--2014 data, specialty and pay revenue would be classified correctly in about \lttvDesignCorrectSpec{} of simulated worlds and distributors' revenue in about \lttvDesignCorrectBDU{}. Canadian programming expenditure, whose historical volatility was \lttvHistSDCPE{} a year, came out at \lttvDesignCorrectCPE{}, so it was set aside in advance as descriptive.

The outcome measure is a single number, \(k\): the observed shortfall below the report's baseline, 2016--2019, as a multiple of the report's own forecast impact, so that \(k=0\) is no effect and \(k=1\) is the forecast. The baseline's error is modelled as a random walk from 2014. The forecasters' own cited inputs (uptake between 10 and 35\%, closures up to two and a half times the assumed rate) define a band of \(k\) that counts as consistent with their forecast. An outcome below the forecast path was to be reported as the forecast exceeded, never as confirmed.

How wrong the baseline was can't be known from the outcomes, so the plan fixed an external check: the report's forecast of US Netflix subscribers, which Canadian regulation couldn't affect, against Netflix's reported US paid memberships \citep{netflix2019form10k}. The report forecast \lttvNetflixFcEighteen{} million for 2018 and Netflix reported \lttvNetflixObsEighteen{} million, which by the pre-stated rule puts the error at \lttvSigmaTech{} a year, below the historical volatility of both revenue series. The CRTC series reproduce the report's 2012--2014 values within 2\% \citep{crtc2017sfs2016, crtc2021sfs2020}.

\subsection{What happened}

Specialty and pay revenue didn't fall. From 2016 to 2019 it ran between \lttvSpecVsBaselineMin{} and \lttvSpecVsBaselineMax{} \lttvSpecAboveBaseline{} the report's no-reform baseline, and between \lttvSpecVsLTTVMin{} and \lttvSpecVsLTTVMax{} \lttvSpecAboveLTTV{} its forecast path. At the pre-stated reading \(k=\lttvKSpec\) (95\% interval \lttvKSpecLo{} to \lttvKSpecHi{}), against a band of \lttvBandSpecLo{} to \lttvBandSpecHi{} implied by the forecasters' own inputs.

Distributors' revenue fell further than forecast. It ran between \lttvBDUVsLTTVMin{} and \lttvBDUVsLTTVMax{} \lttvBDUAboveLTTV{} the report's forecast path, falling from \$\lttvBDURevFifteen{} million in 2015 to \$\lttvBDURevNineteen{} million in 2019. Here \(k=\lttvKBDU\) (\lttvKBDULo{} to \lttvKBDUHi{}), inside the forecasters' band of \lttvBandBDULo{} to \lttvBandBDUHi{}.

Neither verdict is robust to how wrong the baseline is taken to be. Specialty and pay revenue becomes inconclusive once the baseline's error reaches \lttvCutoffSpec{} a year, and distributors' revenue at \lttvCutoffBDU{}. Other defensible readings of the Netflix comparison cross those lines: the 2016 miss gives \lttvSigmaTechSixteen{} a year and a pooled estimate \lttvSigmaTechPooled{}. Part of the 2016 miss is an inherited offset, since the report's 2015 starting figure (\lttvNetflixFcFifteen{} million) was already above Netflix's count (\lttvNetflixObsFifteen{} million) \citep{netflix2018form10k}.

Measured as growth, though, Netflix outran the forecast, by \lttvNetflixObsGrowth{} against \lttvNetflixFcGrowth{} from 2015 to 2018, which puts the error at \lttvSigmaGrowthEighteen{} a year. Across the \lttvNCalibrations{} calibrations I examined, specialty and pay revenue is inconclusive under \lttvNInconclusiveSpec{} and smaller than the forecasters' inputs imply under the rest; distributors' revenue is inconclusive under \lttvNInconclusiveBDU{}. The estimate of \(k\) doesn't move, and no interval for specialty and pay revenue excludes zero \citep{gelman2014statcrisis, gelmanStern2006difference}.

\subsection{The links in the chain}

The report's own inputs can be checked without a counterfactual, and they show where its chain parted from events. Uptake of the entry-level service was small. The CRTC counted \lttvUptakeAprilN{} subscribers five weeks after launch \citep{crtc2016basicpackage} and \lttvUptakeJuneN{} by 30 June 2016 \citep[para.~12]{crtc2016hearing0907}, \lttvUptakeAprilShare{} and \lttvUptakeJuneShare{} of all subscribers, against the report's \lttvByopSixteen{} for 2016. I found no later count.

The modelled pass-through isn't visible in the aggregate. While distributors' revenue fell, what they paid to Canadian channels rose, from \$\lttvPayCanFifteen{} million in 2015 to \$\lttvPayCanNineteen{} million in 2019, and the Canadian share of all affiliation payments stayed near \lttvPayShareFifteen{} \citep[Table~U-T11]{crtc2019cmrbdu}. The aggregates don't show why. Rising fees on a shrinking subscriber base are consistent with it, and so is the CRTC's wholesale code, which the report itself lists as coming into effect in January 2016 \citep[p.~73, n.~82]{nordicity2015television}.

Closures are the one input that came out close to the report. Between \lttvClosVILo{} and \lttvClosVIHi{} of the integrated companies' Category A and B channels stopped operating by 2019--2020, against the report's 10\%. For the independents the bounds run from \lttvClosIndLo{} to \lttvClosIndHi{}, because channels that left licensing for exemption can't be told apart from channels that closed. Canadian programming expenditure, read only descriptively, stayed at or above the baseline path through 2019.

\subsection{What the case shows}

The report's technical claim held in one place and failed in another. At the pre-stated reading, distributors lost revenue on the scale the model predicted, though the entry-level service drew far fewer subscribers than it assumed, so how much of that loss the decisions caused isn't settled here. The chain then broke at the transfer to Canadian channels, whose revenue stayed at or above what the report expected with no reform at all.

Most of the jobs figure put to Parliament comes after that break. Of the \lttvJobsTotal{} FTEs, \lttvFTESpecProd{} (\lttvFTESpecProdShare{}) are in specialty and pay services and independent production, and depend mainly on Canadian channels' revenue and programming spending; the other \lttvFTEBDU{} are distributors' own, and follow from the revenue loss that did occur \citep[Tables~22--23, pp.~91--92]{nordicity2015television}. The public argument went further: the testimony presented an economy-wide, mostly spin-off estimate as media jobs that would be lost, as the direct result of the policy, and quoted a share of programming spending that the report's own figures don't support.

Neither of the authors' later reports for the CRTC that I read returns to these numbers \citep{miller2022rights, nordicity2022harnessing}. The case resembles the forecasts made during the \textit{SFFA} litigation that \textcite{kane2026sffa} checked: the forecasters were interested parties, the forecast was placed before decision-makers, and the outcomes are in public data.
